% Phase-2 fix designs, group G2 (computability / arithmetization), defects R-002..R-006.
% Fragment: plain LaTeX, compilable under amsmath/amssymb when wrapped in a document.
% Line anchors refer to the v1.0 manuscript source/formal-logic-reference.tex (14780 lines).

\subsection*{R-002 --- Bounded \(\mathsf{Comp}\) versus the Cantor-pair numbering (v1.0 lines 7004--7031, 6564--6593)}

\paragraph{Defect.} The theorem at 7015--7028 asserts a bounded \(\mathsf{Comp}(e,x,s)\) recognizing accepting computations in the fixed machine numbering (6567--6585), while 7011--7013 concedes that numbering's iterated Cantor pairing ``comes with no bounded extraction formula and is not the coding used here.''  \(\mathsf{Comp}\) must still decode \(e\) and \(x\) from iterated Cantor codes; the citations (7030--7031) arithmetize other numberings.

\paragraph{Repair.} Option (ii): keep the numbering; make the history coding carry \(\beta\)-extractable \emph{decoding certificates}, so \(\mathsf{Comp}\) becomes bounded.  Replace 7004--7013 by the material below and keep 7015--7027 plus one new clause; 7030--7031 stays.

Lemma A (pairing graph; lands before 7015).  \(\mathsf{P}(x,y,z):\equiv(x+y)\times S(x+y)+(y+y)=z+z\) is quantifier-free, defines the graph of \(\operatorname{pair}_{\mathbb N}\), and \(a,b\le\operatorname{pair}_{\mathbb N}(a,b)\) (\(b\le z\); \(a\le a+b\le(a+b)(a+b+1)/2\le z\)).  Status: complete proof.  Lemma B (G\"odel \(\beta\)).  With the bounded formula
\[
  \mathsf{B}(w,i,v):\equiv\exists t<Sw\,\exists m<Sw\,\exists q<St\,
  \bigl(\mathsf{P}(t,m,w)\land t=d\times q+v\land\exists u<d\,(v+Su=d)\bigr),
\]
writing \(d\) for the term \(S((Si)\times m)\) (\(u\) fresh, \(u\notin\mathrm{Var}(d)\)).  The final conjunct is the \emph{marked-bounded} form of \(v<d\), \(\mathcal N\)-equivalent to the plain abbreviation of 6924--6926 because any witness of \(v+Su=d\) satisfies \(Su\le d\), hence \(u<d\); the marked form is required because the book counts only marked bounded quantifiers as bounded (6936--6937, 6939--6951), while a bare comparison \(r<t\) expands to an unbounded existential.  Claims: (a) for all \(w,i\) exactly one \(v\le w\) satisfies \(\mathsf{B}\) in \(\mathcal N\); write \(\beta(w,i)=v\); (b) every finite sequence \(a_0,\dots,a_{k-1}\) is \(\beta(w,\cdot)\upharpoonright k\) for some \(w\).  Proof of (a): \(\operatorname{pair}_{\mathbb N}\) is a bijection, so \(w\) determines \(t,m\) uniquely, with \(t,m\le w\) by Lemma A, matching the displayed bounds; Euclidean division of \(t\) by \(d\ge1\) determines \(q,v\) uniquely, with \(q\le t<St\) and \(v\le t\le w\).  Proof of (b): put \(m:=N!\) with \(N:=\max(k,a_0,\dots,a_{k-1})+1\); the moduli \(1+(i+1)m\) are pairwise coprime (a common prime divides \((j-i)m\) but not \(m\), hence divides \(j-i\); yet \(0<|j-i|<k\le N\) forces it to divide \(m\)); each \(a_i<m\); Chinese remaindering gives \(t\equiv a_i\pmod{1+(i+1)m}\); put \(w:=\operatorname{pair}_{\mathbb N}(t,m)\).  G\"odel 1931 and Boolos--Burgess--Jeffrey 2007, ch.~16, are parallel reading.  Status: complete proof.

Definition (history coding; replaces 7004--7013).  \(s\) is a \emph{history code} for \((e,x)\) iff, under \(\beta(s,\cdot)\) with a fixed index scheme (families at \(\operatorname{pair}_{\mathbb N}(t,j)\)), it contains: (h1) the binary digits of \(x\) with certified expansion \(v_0=d_0\), \(v_{j+1}=2v_j+d_{j+1}\), \(v_{L-1}=x\), \(v_j\le x\), \(d_j\le1\); (h2) the table of \(e\): a chain \(\mathsf P(k,c_0,e)\), \(\mathsf P(r_i,c_{i+1},c_i)\) (\(i<k-1\)), \(r_{k-1}=c_{k-1}\), rows unpacked by fixed-depth \(\mathsf P\)-steps, plus the validity checks of 6567--6576 under the designated-slot normalization of the R-005 repair; (h3) configurations \(C_0,\dots,C_T\), tapes coded in the tape-alphabet base with certified expansions, the first from (h1), each step matching the decoded table by bounded digit comparisons; (h4) final state accepting, tape prefix before the first blank a valid binary numeral (6327--6338).  Every auxiliary value is an entry of \(s\) (Lemma B(a)) or a Cantor-code component (Lemma A), so all quantifiers are bounded by terms in \(e,x,s\); and every comparison conjunct \(r<t\) or \(r\leq t\) occurring in \(\mathsf{Comp}\) is written in its marked-bounded form \(\exists u<t\,(r+Su=t)\), respectively \(\exists u\leq t\,(u+r=t)\) with \(\leq\) as registered by the R-003 repair (\(u\) fresh), each \(\mathcal N\)-equivalent to the plain abbreviation, so the matrix is bounded under the counting rule of 6936.  Repaired theorem: \(\mathsf{Comp}(e,x,s)\) is bounded; \(\mathcal N\models\mathsf{Comp}(\overline e,\overline x,\overline s)\) iff \(s\) is a history code for \((e,x)\); and \(\varphi_e(x)\) is defined iff \(\mathcal N\models\exists s\,\mathsf{Comp}(\overline e,\overline x,s)\) (existence via Lemma B(b); invalid \(e\) fail (h2), matching divergence).  Status: theorem proof sketch with plan (h1)--(h4); Lemmas A, B complete proof.

Finally, add at 7449--7454 one sentence noting that the coded-computation-sequence apparatus behind \cref{thm:representability}'s \(\Sigma^0_1\) clause (the construction paragraph at 7351--7357) is the \(\beta\) coding of Lemma B, so a single in-document coding serves machine tables, histories, and proof sequences; this discharges the register's one-coding acceptance for \(\mathsf{Prov}_T\).

\paragraph{Acceptance.} (1) 7011--7013 deleted; history coding in-document.  (2) Every displayed quantifier bounded by a term in \(e,x,s\); every comparison conjunct in marked-bounded form; monotonicity proved.  (3) 7040--7051, 7068--7074, 13735--13760 type-check unchanged.  (4) The \(\Sigma^0_1\) presentations of \(\mathsf{Comp}\) and of \(\mathsf{Prov}_T\) (7440--7454) both point at Lemma B's apparatus.

\subsection*{R-003 --- The two \(\mathsf{Q}\)-lemmas for Rosser's bounded comparison (v1.0 lines 7663--7758, 6924--6937, 7210--7222)}

\paragraph{Defect.} The Rosser proof plan (7682--7700) rests on ``evaluating the bounded quantifier in \(\mathsf Q\)'' and on the exhibited proof witnessing ``the bounded existential for every larger candidate.''  Neither is backed: \(\leq\) is never introduced (only \(<\), 6924--6926), and the two needed \(\mathsf Q\)-facts --- bounded-case enumeration and dichotomy --- are neither stated nor cited.

\paragraph{Repair.} (1) Register, next to \(<\) at 6924--6937: \(r\leq t:\equiv\exists z\,(z+r=t)\), \(z\) least fresh; \(\exists x\!\leq\! t\,\varphi:\equiv\exists x(x\leq t\land\varphi)\), \(\forall x\!\leq\! t\,\varphi:\equiv\forall x(x\leq t\to\varphi)\) for \(x\notin\mathrm{Var}(t)\), counted as bounded, like \(<\).  Record: \(\mathcal N\models x\leq t\leftrightarrow x<St\), \emph{not} claimed in \(\mathsf Q\); the witness sits \emph{left} of \(+\) so the lemmas need only \(\mathsf Q\)'s right-recursion axioms.

(2) Two numbered lemmas after \(\mathsf Q\)'s definition (7222), proved by metainduction on \(n\); each instance is a \(\mathsf Q\)-derivation from \(Sx\neq0\), \(S\)-injectivity, \(x\neq0\to\exists y\,x=Sy\), \(x+0=x\), \(x+Sy=S(x+y)\).

Lemma Q1 (bounded-case enumeration) and Lemma Q2 (dichotomy).  For every \(n\in\mathbb N\):
\[
  \mathsf Q\vdash\forall y\,\bigl(y\leq\overline n\to(y=\overline0\lor y=\overline1\lor\dots\lor y=\overline n)\bigr),
  \qquad
  \mathsf Q\vdash\forall y\,\bigl(y\leq\overline n\lor\overline n\leq y\bigr).
\]
Proof of Q1: \(n=0\): from \(z+y=0\), if \(y=Sw\) then \(S(z+w)=z+Sw=0\) contradicts \(Sx\neq0\); so \(y=0\).  Step: from \(z+y=S\overline n\): \(y=0\) gives the first disjunct; \(y=Sw\) gives \(S(z+w)=S\overline n\), \(S\)-injectivity gives \(z+w=\overline n\), so \(w\leq\overline n\); the induction hypothesis gives \(w=\overline0\lor\dots\lor w=\overline n\); apply \(S\).  Proof of Q2: \(n=0\): \(\overline0\leq y\), witness \(z:=y\), by \(y+0=y\).  Step: fix \(y\); if \(y=0\), \(y\leq S\overline n\) (witness \(S\overline n\)); if \(y=Sw\), the hypothesis on \(w\) gives either \(z+w=\overline n\), whence \(z+Sw=S(z+w)=S\overline n\) and \(y\leq S\overline n\), or \(u+\overline n=w\), whence \(u+S\overline n=S(u+\overline n)=Sw=y\) and \(S\overline n\leq y\).  Status of both: complete proof.

(3) Rewrite 7682--7700 to invoke the lemmas by label, with fixed point \(\mathsf Q\vdash R_T\leftrightarrow\forall y\,(\mathsf{Proof}_T(y,\overline{\ulcorner R_T\urcorner})\to\exists z\!\leq\! y\,\mathsf{Proof}_T(z,\overline{\ulcorner\neg R_T\urcorner}))\).  Branch \(T\vdash R_T\), least proof code \(p\): representability gives \(T\vdash\mathsf{Proof}_T(\overline p,\overline{\ulcorner R_T\urcorner})\), hence \(T\vdash\exists z\!\leq\!\overline p\,\mathsf{Proof}_T(z,\overline{\ulcorner\neg R_T\urcorner})\); consistency gives \(T\nvdash\neg R_T\), so numeralwise \(T\vdash\neg\mathsf{Proof}_T(\overline q,\overline{\ulcorner\neg R_T\urcorner})\) for \(q\leq p\); Lemma Q1 turns the bounded existential into a finite disjunction: inconsistency.  Branch \(T\vdash\neg R_T\), least code \(q\): reason in \(T\) from \(\mathsf{Proof}_T(y,\overline{\ulcorner R_T\urcorner})\); Lemma Q2 gives \(y\leq\overline q\lor\overline q\leq y\); first case: Lemma Q1 plus \(T\vdash\neg\mathsf{Proof}_T(\overline k,\overline{\ulcorner R_T\urcorner})\) for \(k\leq q\) (from \(T\nvdash R_T\)) yields a contradiction; second case: witness \(z:=\overline q\).  Hence \(T\vdash R_T\), contradicting consistency.  Status: Rosser stays proof sketch, now citing Q1/Q2 by label; the solution at 13830--13838 names both lemmas.

\paragraph{Acceptance.} (1) \(\leq\) is registered with its expansion, the bounded-counting rule, and the non-claim of \(\mathsf Q\)-equivalence with \(x<St\); no arithmetic object-formula uses \(\leq\) earlier.  (2) Q1 and Q2 appear exactly as displayed, proved by metainduction from the \(\mathsf Q\)-axioms, labeled complete proof.  (3) Each Rosser branch names which lemma discharges which step.

\subsection*{R-004 --- Jump definition and the partiality convention (v1.0 lines 7134--7189, 6576--6582)}

\paragraph{Defect.} The jump (7154--7161) uses ``the \(e\)-th oracle machine with oracle \(B\) \emph{halts} on input \(e\),'' while \(K\) (6678--6681) uses ``\(\varphi_e(e)\) is defined,'' which demands an accepting halt with valid numerical output (6576--6582); halting also covers rejection and invalid output.  So the claim at 7163--7165 that the \(K\) diagonal ``relativizes'' fails literally and is unlabeled prose; 7142--7144 also never fixes the oracle I/O convention.

\paragraph{Repair.} (1) Replace 7154--7161.  Fix an effective numbering of oracle machines analogous to 6567--6585 (same normalization, query rows included, invalid codes diverge everywhere).  For valid \(e\) and \(B\subseteq\mathbb N\), let \(\varphi^B_e(x)\) be the numerical output when the decoded machine with oracle \(B\), on \(\mathrm{bin}(x)\), halts accepting with a valid numerical output, and undefined when it rejects, diverges, or accepts without one --- 6576--6582 verbatim, relativized.  Define \(B':=\{e\in\mathbb N:\varphi^B_e(e)\text{ is defined}\}\), \(B^{(0)}:=B\), \(B^{(n+1)}:=(B^{(n)})'\).  Amend 7142--7145: oracle machines ``compute'' by this discipline; ``computably enumerable relative to \(B\)'' means: domain of some \(\varphi^B_e\).

(2) New lemma (relativized adequacy).  Every partial function computed with oracle \(B\) occurs as \(\varphi^B_e\); a relativized universal machine gives \(U^B(e,x)\simeq\varphi^B_e(x)\); and a relativized parameter theorem holds: for every \(m,n\geq1\) there is a total computable \(\mathsf{par}_n^m\) with \(\varphi^B_{\mathsf{par}_n^m(e,\bar a)}(\bar x)\simeq\varphi^B_e(\bar a,\bar x)\) for every oracle \(B\) simultaneously.  Route: R-005's adequacy proof (renaming commutes with the step relation, queries included), the universal simulation of 6608--6618 with queries passed through, and the compiler construction of 6632--6641 with query rows copied unchanged.  Status: proof sketch.

(3) Promote 7163--7165 to a numbered theorem.  Jump theorem: for every \(B\subseteq\mathbb N\), (a) \(B'\) is computably enumerable relative to \(B\); (b) \(B'\not\leq_T B\).  Proof (relativized diagonal): (a) \(B'\) is the domain of \(e\mapsto U^B(e,e)\).  (b) If an oracle machine with oracle \(B\) decides \(B'\) via total \(h^B\), define \(d^B(e):=0\) if \(h^B(e)=0\), undefined otherwise; \(d^B\) is oracle-computable (run the decider; write \(0\) and accept, or loop); fix \(b\) with \(\varphi^B_b\simeq d^B\) by the adequacy lemma; then \(b\in B'\) iff \(\varphi^B_b(b)\) is defined iff \(h^B(b)=0\) iff \(b\notin B'\).  Status: proof sketch --- the displayed diagonal calculation is complete, but clause (a) rests on the relativized universal machine and clause (b) on the occurrence clause of the relativized adequacy lemma, both proof sketches; the status policy (313--319) and the release-checklist rule (14346) forbid complete-proof status over a sketch dependency.

(4) New remark after 7189, keeping citations sound.  (i) \(K\equiv_m\varnothing'\): an ordinary table embeds as a query-free oracle table; an oracle-\(\varnothing\) machine is simulated by answering queries negatively.  (ii) The halting-based jump of the cited sources is \(\equiv_m\) this \(B'\): \(e\mapsto e^*\) simulating \(e\) and, on any halt, writing \(0\) and accepting (halting \(\leq_m\) defined); \(e\mapsto e^{**}\) diverging unless the halt is accepting with valid output (defined \(\leq_m\) halting).  All four maps are computable table transformations, uniform via the parameter theorem (\cref{thm:parameter}).  Jump monotonicity --- \(A\leq_T B\) implies \(A'\leq_m B'\) --- is stated and used explicitly: replace each \(A\)-query of an oracle-\(A\) machine by a run of a fixed oracle-\(B\) decider for \(A\), then align self-application indices by the relativized parameter theorem of the adequacy lemma with padding.  Hence Post's theorem (7167--7175) is insensitive to the convention; the citations stand.  Status: proof sketch --- the four table transformations are elementary, but the monotonicity step runs through the relativized parameter theorem, part of the sketch-status adequacy lemma.

\paragraph{Acceptance.} (1) \(B'\) is defined via ``\(\varphi^B_e(e)\) is defined,'' three-way undefinedness relativized verbatim; 7142--7145 fixes the oracle I/O convention.  (2) The jump theorem is a numbered statused result argued by the relativized diagonal invoking the labeled adequacy lemma, with a status label matching its sketch dependencies.  (3) The remark states \(K\equiv_m\varnothing'\), the \(\equiv_m\)-bridge to the halting convention, and the jump-monotonicity step it uses, each with honest status.

\subsection*{R-005 --- The enumeration claim for the fixed numbering (v1.0 lines 6564--6593, decisive at 6582--6585)}

\paragraph{Defect.} The claim that ``every unary partial computable function occurs in this numbering'' and, via the tuple convention, every \(k\)-ary one (6582--6585), is unlabeled prose with no proof, status, or citation.  Every construction that picks an index for an informally described machine (6692--6699, 6724, 6795--6816, 13735--13760) depends on it, and a bare citation cannot discharge it: the numbering is bespoke (normalization, validity constraint \(\{0,1\}\subseteq\mathcal A_{\mathrm{in}}\), wrapper).

\paragraph{Repair.} (1) Amend 6567--6569 to fix the normalization completely: a description is \emph{normalized} iff its tape symbols are the numbers \(0,\dots,|\mathcal A_{\mathrm{tape}}|-1\) with the designated assignments \(\sqcup\mapsto0\), digit \(0\mapsto1\), digit \(1\mapsto2\), and its states are the numbers \(0,\dots,|Q|-1\) with \(q_0\mapsto0\), \(q_{\mathrm{acc}}\mapsto1\), \(q_{\mathrm{rej}}\mapsto2\); the remaining numbers are assigned arbitrarily, and the row list is the lexicographic listing by numbered key \((q,a)\) over the non-halting domain.  (An ``order of first appearance'' clause would be circular: the lexicographic row order presupposes the numbering.)  No canonical machine-to-code map is needed, since indices are not unique (6587); valid codes thus have pairwise distinct designated states, and machines whose start state is halting (permitted by 6233--6234) are captured up to behavior by the surrogate-start step in the theorem's proof below.  Without the designated slots ``valid code'' is not well defined, and the validity clause \(\{0,1\}\subseteq\mathcal A_{\mathrm{in}}\) together with the input and output conventions (6296--6338) is to be read at the designated digit numbers \(1,2\).

(2) New lemma (renaming invariance; lands after 6593).  If \(\iota=(\iota_Q,\iota_A)\) bijectively renames states and tape symbols preserving the roles \(q_0,q_{\mathrm{acc}},q_{\mathrm{rej}},\sqcup,0,1\), then \(C\to_{\mathcal T}C'\) iff \(C^\iota\to_{\mathcal T^\iota}(C')^\iota\); hence \(\mathcal T\) and \(\mathcal T^\iota\) accept, reject, or diverge together on every input, with corresponding tapes.  Proof: one-step induction from 6248--6294.  Status: complete proof.

(3) Replace 6582--6585 by a numbered theorem.  Adequacy of the numbering (enumeration theorem): (a) every \(\varphi_e\) is partial computable; (b) for every \(k\geq1\) and partial computable \(f:\mathbb N^k\rightharpoonup\mathbb N\) there is \(e\) with \(\varphi_e(\bar x)\simeq f(\bar x)\) for all \(\bar x\in\mathbb N^k\), under the arity-\(k\) convention of 6588--6593.  Proof route: (a) is the wrapper argument at 6579--6582, folded in.  (b) Let \(\mathcal T\) witness \(f\) per 6340--6349.  First, \(\{0,1\}\subseteq\mathcal A_{\mathrm{in}}\) automatically: \(\operatorname{code}_k\) is surjective for every \(k\geq1\) (iterated pairing bijections), so the inputs \(\mathrm{bin}(\operatorname{code}_k(\bar x))\) exhaust all binary numerals; initial configurations exist only for words in \(\mathcal A_{\mathrm{in}}^*\) (6280--6288), and 6340--6349 demands an accepting or infinite computation on each such input; so both digits lie in \(\mathcal A_{\mathrm{in}}\), and no alphabet extension is needed.  Second, since 6233--6234 requires only \(q_{\mathrm{acc}}\neq q_{\mathrm{rej}}\), the start state may coincide with a halting state; in that case first replace it by a fresh start state \(q_0'\) with \(\delta(q_0',a)=(q_0,a,R)\) for every tape symbol \(a\) --- the one extra step leaves the tape unchanged, and the output convention ignores the final head position and off-prefix cells (6331--6338), so the accepted outputs and divergence behavior are unchanged.  Then the three designated states and three designated symbols are pairwise distinct; choose any role-preserving bijective renaming onto the designated-slot numbering, obtaining a normalized \(\mathcal T^*\) that behaves identically by the renaming lemma; list its rows lexicographically by numbered key (the row list is the full graph of \(\delta\), total on non-halting pairs by 6220--6246), encode as \(\operatorname{pair}_{\mathbb N}(k,\operatorname{code}_k(\text{rows}))\); the code \(e\) is valid by construction, and by the lemma \(e\)'s behavior on every \(\mathrm{bin}(\operatorname{code}_k(\bar x))\) is \(\mathcal T\)'s, i.e.\ \(\varphi_e(\bar x)\simeq f(\bar x)\).  Retain the Kleene 1952 citation (enumeration and normal form over standard numberings) explicitly as parallel reading; the in-document proof is primary because the numbering is bespoke.  Status afterwards: complete proof.

\paragraph{Acceptance.} (1) No unlabeled adequacy prose remains at 6582--6585; the claim exists only as the numbered theorem, clause (b) quantifying over all \(k\geq1\) with \(\varphi_e(\bar x)\simeq f(\bar x)\).  (2) The normalization convention (designated numbers for \(\sqcup,0,1,q_0,q_{\mathrm{acc}},q_{\mathrm{rej}}\)) is fixed in the numbering paragraph and used by the proof.  (3) The diagonal construction (6692--6699) and the Rice proof (6803--6816) cite the theorem by label when picking indices.

\subsection*{R-006 --- Closure-class versus prefix hierarchy (v1.0 lines 6939--7002, 7030--7189)}

\paragraph{Defect.} The hierarchy definition (6952--6971) builds \(\Sigma^0_{n+1}\), \(\Pi^0_{n+1}\) as \emph{closure classes} (under \(\land,\lor\), both bounded quantifiers, one unbounded quantifier), while every hierarchy fact discharged by citation (7030--7031, 7078--7080, 7082--7097, 7167--7189) concerns the literature's \emph{prefix} hierarchy.  Nothing identifies the two, so the citations do not apply; the missing ingredient is bounded-quantifier exchange.

\paragraph{Repair.} Insert after 7002:

Lemma C (contraction).  Let \(\mathsf P(x,y,z)\) be the quantifier-free pairing-graph formula of the R-002 repair's Lemma A, with its monotonicity bound \(a,b\le\operatorname{pair}_{\mathbb N}(a,b)\); the two designs land at adjacent sites (before 7015, after 7002), so the manuscript states \(\mathsf P\) and its lemma exactly once and Lemma C cites it.  Then, for \(z\) fresh (\(z\) not free in \(\theta\), \(z\notin\{x,y\}\)): \(\mathcal N\models\exists x\exists y\,\theta\leftrightarrow\exists z\,\exists x\!<\!Sz\,\exists y\!<\!Sz\,(\mathsf P(x,y,z)\land\theta)\), dually for \(\forall\forall\) with \(\to\).  Status: complete proof (via Lemma A).

Lemma D (bounded-quantifier exchange).  For every formula \(\theta\) and term \(t\) with \(y,w\notin\mathrm{Var}(t)\), \(w\) fresh:
\[
  \mathcal N\models\forall y\!<\!t\,\exists x\,\theta\leftrightarrow\exists w\,\forall y\!<\!t\,\exists x\!<\!w\,\theta,
  \quad
  \mathcal N\models\exists y\!<\!t\,\forall x\,\theta\leftrightarrow\forall w\,\exists y\!<\!t\,\forall x\!<\!w\,\theta .
\]
Proof: right-to-left immediate; left-to-right: finitely many \(y<\mathrm{val}(t)\), choose witnesses \(x_y\), put \(w:=\max_y x_y+1\); the second display is dual.  Equivalences \emph{in \(\mathcal N\)}, for boldface classes only; no \(\mathsf Q\)- or \(\mathsf{PA}\)-provability is claimed.  Status: complete proof.

Theorem E (prefix normal form).  For every \(n\geq0\), relations defined by prefix-form \(\Sigma^0_n\) formulas (\(n\) alternating unbounded blocks starting \(\exists\), bounded matrix) are closed, up to \(\mathcal N\)-equivalence of the defining formulas, under \(\land\), \(\lor\), bounded \(\exists\), bounded \(\forall\), unbounded \(\exists\); dually for \(\Pi^0_n\).  Hence, by induction on the clauses of 6952--6971, every closure-class formula is \(\mathcal N\)-equivalent to a prefix-form formula of the same class; since conversely every prefix-form formula is generated by those clauses, the boldface classes of 6974--7002 coincide with the cited sources'.  Proof route (execute literally): induction on \(n\), both classes at once, \(\Pi\) dual.  \(n=0\): all operations except unbounded \(\exists\) stay bounded syntactically.  \(n\geq1\), relation \(\exists x\,B\), \(B\) prefix-\(\Pi^0_{n-1}\): (i) \(\land,\lor\): rename, pull existentials out, contract by Lemma C; the residual bounded \(\exists\)'s land on prefix-\(\Pi^0_{n-1}\) formulas, absorbed at level \(n-1\).  (ii) bounded \(\exists\): \(\exists y\!<\!t\,\exists x\,B\equiv\exists x\,\exists y\!<\!t\,B\) by pure logic; absorb at level \(n-1\).  (iii) bounded \(\forall\) (collection): by Lemma D, \(\forall y\!<\!t\,\exists x\,B\) is \(\mathcal N\)-equivalent to \(\exists w\,\forall y\!<\!t\,\exists x\!<\!w\,B\); absorb both bounded quantifiers at level \(n-1\), prepend \(\exists w\).  (iv) unbounded \(\exists\): contract by Lemma C.  Base: a level-\((n-1)\) formula converts by hypothesis and is padded with dummy quantifiers.  Status afterwards: complete proof.  Cite Theorem E at 7078--7080 for the \(\mathsf{TOT}\) bound and at 7030--7031, 7097, 7178--7189 as license for the prefix-based citations.

\paragraph{Acceptance.} (1) Lemma D appears exactly as displayed, proved by the finite-witness maximum, \(\mathcal N\)-equivalence only.  (2) Theorem E's induction covers all five operations for both classes at every level, bounded-\(\forall\)-over-\(\Sigma\) via Lemma D, contractions via Lemma C.  (3) All four citation sites reference Theorem E.
